% Self-contained supplement to C.1.  No preamble or bibliography command.
\paragraph{C.1g (Proved: a specified generic quintic and an infinite-order difference of lines).}
We now give the missing line argument, with a specified function field.
This proof uses the local period-map formula and the elementary cyclic
projection idea of \citet[Section III]{cyCeresaCollino1983}; the polynomial
interpolation, the non-torsion specialization, and the passage to a
specified geometric generic fibre are given explicitly below. We do not
claim that every complex specialization has the asserted property.

Let
\[
 \mathcal M=\{e\in\mathbb Z_{\ge0}^{5}:|e|=5,\quad
                  e_2+e_3+e_4>0,\quad e_0+e_1+e_4>0\},
 \qquad e_*=(0,0,0,0,5).
\]
There are $114$ monomials in $\mathcal M$. For algebraically independent
indeterminates $t_e$, $e\in\mathcal M\setminus\{e_*\}$, set
\[
 K=\mathbb Q(t_e:e\in\mathcal M\setminus\{e_*\}),\qquad
 F=x_4^5+\sum_{e\ne e_*}t_e x^e,
 \qquad V=\{F=0\}\subset\mathbb P^4_K.
\]
The two specified lines are
\[
 L_0=\{x_2=x_3=x_4=0\},\qquad
 L_1=\{x_0=x_1=x_4=0\}.
\]
Then $V$ is smooth, $[L_0]-[L_1]$ is homologically trivial, and
\begin{equation}\label{eq:specified-quintic-lines}
 m([L_0]-[L_1])\ne0
 \quad\hbox{in }\mathrm{CH}^2(V_{\overline K})/
                    \mathrm{CH}^2(V_{\overline K})_{\mathrm{alg}}
 \quad\text{for every }m\ne0.
\end{equation}
The same assertion holds after every embedding $\overline K\hookrightarrow
\mathbb C$. It therefore supplies both a field and an equation for the
counterexample, without making an assumption about fields of definition
of analytic exceptional loci.

\emph{The parameter space and all its infinitesimal deformations.}
Let $T$ be the nonempty smooth-quintic locus in the affine parameter
space just displayed. The complete linear system of quintics containing
$L_0\cup L_1$ has that union as its base locus. Bertini gives smoothness
away from the lines, and its normal derivatives generate the conormal
space at every point of either line. Thus a general member is smooth
there as well. This proves that $T$ is nonempty and the generic $V$ is
smooth.

We will also need that this family supplies every abstract first-order
deformation at a general point. Here is a direct check. In the same
linear system consider the polynomial
\[
 F_* =x_2x_0^4+x_3x_0^2x_1^2+x_4x_1^4
       +x_0x_2^4+x_1x_2^2x_3^2+x_4x_3^4+x_4^5.
\]
On each of the two lines, its three normal derivatives are
$(u^4,u^2v^2,v^4)$. The resulting map
\[
 H^0(\mathbb P^1,\mathcal O(1))^{\oplus3}
       \longrightarrow H^0(\mathbb P^1,\mathcal O(5))
\]
is an isomorphism: multiplication sends the six standard basis elements
to
\[
 u^5,\quad u^4v,\quad u^3v^2,\quad u^2v^3,\quad uv^4,\quad v^5.
\]
This is an open condition on
the coefficients. We use $F_*$ only to certify that open condition,
and do not need it itself to be a smooth quintic. Intersecting this
nonempty open set with the smooth locus gives smooth quintics for which
$H^0(N_{L_j/V})=H^1(N_{L_j/V})=0$ for both $j$.

The incidence space of a quintic with two ordered skew lines is a
projective bundle over the smooth space of ordered skew pairs. At the
points just described its map to the full quintic parameter space has
invertible differential: every first-order perturbation of the equation
has a unique compatible first-order perturbation of each line, by the
six-by-six isomorphism. The incidence map is consequently dominant.
The group $\mathrm{PGL}_5$ acts transitively on ordered skew pairs,
and its directions induce only trivial abstract deformations. It follows
that the Kodaira--Spencer map from $T$ onto $H^1(V,T_V)$ is surjective
on a nonempty open subset $T^\circ$.

For clarity, all embedded deformations do give all abstract ones here.
The normal sequence and the Euler sequence give
\[
 H^0(V,\mathcal O_V(5))\longrightarrow H^1(V,T_V)\longrightarrow
 H^1(V,T_{\mathbb P^4}|_V)=0.
\]
The last vanishing follows from $H^1(V,\mathcal O_V(1))=0$ and
$H^2(V,\mathcal O_V)=0$, themselves obtained from the hypersurface
sequence and the cohomology of line bundles on $\mathbb P^4$.

\emph{The inversion assertion needed for this family.}
Adjunction gives $K_V\simeq\mathcal O_V$. A nonzero holomorphic
three-form $\Omega$ is nowhere zero, so contraction gives an isomorphism
of sheaves $T_V\simeq\Omega_V^2$ and hence
\[
 H^1(V,T_V)\xrightarrow{\sim}H^{2,1}(V).
\]
On a simply connected small ball in $T^\circ$, trivialize the integral
homology local system and choose a holomorphic generator $\Omega_t$.
For any nonzero integral class $\alpha$ put
$f_\alpha(t)=\int_{\alpha_t}\Omega_t$. This holomorphic function is
not identically zero. Indeed, if it vanished identically, its first
derivatives would vanish. The period-map differential is contraction
of the Kodaira--Spencer class with $\Omega$
\citep[Theorem (1.27), p.~816]{cyGriffiths1968}. Surjectivity just proved
would therefore make $\alpha$ annihilate $H^{3,0}\oplus H^{2,1}$.
Since $\alpha$ is real, it would also annihilate the conjugate space;
Poincar\'e duality would imply $\alpha=0$, a contradiction.

On each such ball there are countably many integral classes.
Outside their proper analytic zero sets, therefore,
\begin{equation}\label{eq:quintic-no-period-kernel}
 \int_\alpha\Omega=0,\quad \alpha\in H_3(V,\mathbb Z)
                     \quad\Longrightarrow\quad \alpha=0.
\end{equation}
We use here that $H_3(V,\mathbb Z)$ is torsion-free. The integral
Lefschetz theorem gives $H_2(V,\mathbb Z)=\mathbb Z$;
the universal coefficient theorem and Poincar\'e duality then give the
claimed torsion-freeness in degree three.

For any algebraic family of curves on a fixed such $V$, parameterized
by a smooth projective curve $D$, the cylinder map
$H_1(D,\mathbb Z)\to H_3(V,\mathbb Z)$ pairs trivially with $\Omega$.
Its swept surface has complex dimension two, so the pulled-back
holomorphic three-form vanishes. By
\eqref{eq:quintic-no-period-kernel} this cylinder map is zero.
The Abel--Jacobi map of the family is holomorphic and its map on
$H_1$ is this cylinder map; this also follows directly by integrating
over the swept chains. A map from a smooth projective curve to a
complex torus which is zero on $H_1$ is constant: its pullbacks of
holomorphic one-forms have zero periods and hence are zero, so its
differential is zero. Differences of fibres thus have Abel--Jacobi
class zero. Closing parameter curves and using refined fibres treats
every generator of algebraic equivalence. We have proved
\[
 AJ\bigl(\mathrm{CH}^2(V)_{\mathrm{alg}}\bigr)=0
\]
for this analytic very general set. The holomorphicity used here is
the construction in \citet[Theorem (2.20), pp.~826--827]{cyGriffiths1968}.

\emph{One smooth fibre where the chosen normal function is not torsion.}
We give the data, including the auxiliary surface. In the plane
$H=\{x_0=0\}\subset\mathbb P^3$ take the smooth Fermat curve
\[
 C:\ x_1^5+x_2^5+x_3^5=0.
\]
Its genus is six. Over the function field of $C\times C$, write its
two generic points in the chart $x_3=1$ as
$z_a=(0,a,b,1)$ and $z_b=(0,c,d,1)$, so
$a^5+b^5+1=c^5+d^5+1=0$ and $b\ne d$.
For every $m>0$, $m[z_a-z_b]\ne0$ in $\operatorname{Pic}^0(C)$:
the map $(z_a,z_b)\mapsto[m]([z_a]-[z_b])$ is not the zero morphism,
since the Abel map is nonconstant and multiplication by $m$ on its
Jacobian is finite. Its value at the generic point cannot be zero.

Take the two skew lines in $\mathbb P^3$
\[
 \ell_a(t,s)=(t,as,bs,s),\qquad
 \ell_b(t,s)=(t,t+cs,ds,s).
\]
Their intersections with $H$ are $z_a,z_b$; they are skew because
$b\ne d$. Define linear forms and a quartic polynomial by
\[
 S=\frac{x_2-bx_3}{d-b},\qquad T=x_1-ax_3-(c-a)S.
\]
\[
 Q(T,S)=T^4+5cT^3S+10c^2T^2S^2+10c^3TS^3+5c^4S^4.
\]
Then $T=S=0$ on $\ell_a$ and $(T,S)=(t,s)$ on $\ell_b$.
Put
\[
 K_0=x_1^5+x_2^5+x_3^5-x_0Q(T,S),\qquad
 K=K_0+x_0G,
\]
where $G$ is a generic element of
$H^0(\mathbb P^3,\mathcal I_{\ell_a\cup\ell_b}(4))$.
Thus $K$ contains both lines and $K|_H$ is exactly the fixed equation
of $C$. These identities follow by substitution, using
$tQ(t,s)=(t+cs)^5-c^5s^5$.

We justify smoothness of $\{K=0\}$, including the base locus.
The projective linear system spanned by $K_0$ and
$x_0H^0(\mathcal I_{\ell_a\cup\ell_b}(4))$ has base locus supported on
$C\cup\ell_a\cup\ell_b$. Bertini applies off it. At every point of
$C$, the derivative within $H$ of its fixed smooth equation is
nonzero. At a point of either line outside $H$, the variations $x_0G$
generate both normal derivatives: products of one linear equation of
each skew line generate the conormal of either line along that line,
and we multiply them by degree-two polynomials nonzero at the point.
Singularity therefore imposes two independent linear conditions at
a point of a one-dimensional line. The union of these bad conditions
has positive codimension in the parameter space. A generic $G$ avoids
it, and $K$ is smooth everywhere.

Now set
\[
 W:\ x_0x_4^4+K(x_0,x_1,x_2,x_3)=0\quad\subset\mathbb P^4,
 \qquad p=(0,0,0,0,1).
\]
This threefold is smooth. A hypothetical singular point has
$x_0x_4^3=0$. When $x_4=0$, smoothness of $K$ excludes it.
When $x_0=0$ and $(x_1,x_2,x_3)\ne0$, smoothness of $C$ excludes
it. The remaining point is $p$, where the $x_0$ derivative is one.
The two lines $\ell_a,\ell_b$ lie on $W$ in $x_4=0$ and avoid $p$.

Let $W'=\operatorname{Bl}_pW$ and $B=\operatorname{Bl}_C\mathbb P^3$.
Projection from $p$ lifts to a morphism $g:W'\to B$. Here is the
needed verification of this lift. In the blow-up of $\mathbb P^4$
at $p$, write its points as $[tu_0:tu_1:tu_2:tu_3:v]$ over
$[u_0:u_1:u_2:u_3]\in\mathbb P^3$. The strict transform has equation
\[
 u_0v^4+t^4K(u)=0.
\]
The ideal $(u_0,K)$ of $C$ pulls back to $(u_0)$ on $t\ne0$, and
to $(K)$ on $v\ne0$. It is invertible, so the universal property
of the blow-up gives $g$. Each proper transform of $\ell_a,\ell_b$
maps isomorphically to the corresponding proper transform in $B$.

Write $j:E\hookrightarrow B$ and $q:E\to C$ for the exceptional
divisor and its projection. The inverse blow-up formula in C.1f
has second component $-q_*j^*:\mathrm{CH}^2(B)\to\mathrm{CH}^1(C)$.
The proper transform of either line meets $E$ transversally once,
over its point $z_a$ or $z_b$. Consequently
\[
 -q_*j^*g_*([\ell_a']-[\ell_b'])=-[z_a-z_b].
\]
The blow-up of a point does not change $H^3$ or the intermediate
Jacobian, while the blow-up of $\mathbb P^3$ in $C$ identifies its
intermediate Jacobian with $\operatorname{Pic}^0(C)$. Abel--Jacobi
functoriality, equivalently integration along this correspondence,
therefore sends $AJ_W(\ell_a-\ell_b)$ to $-[z_a-z_b]$.
The latter is not torsion, so the former is not torsion.

This is a point of the specified family $T$ after an explicit change
of coordinates. Use as the five column vectors
\[
 e_0,\ z_a,\ e_0+e_1,\ z_b,\ q_0=(1,0,0,0,1),
\]
where the first four have fifth coordinate zero. Their determinant
is $b-d$ up to sign. Also $W(q_0)=1$, because $K$ vanishes at
$(1,0,0,0)$ on $\ell_a$. This matrix takes the standard ordered
pair $L_0,L_1$ to $\ell_a,\ell_b$, and the coefficient of $x_4^5$
in the transformed equation is one. All coefficients can be
embedded in $\mathbb C$, preserving the generic points and the
non-torsion conclusion.

\emph{Completion over the specified function field.}
Integral Lefschetz and duality identify $H^4(V,\mathbb Z)$ with
$\mathbb Z$ via degree of a curve, so $L_0-L_1$ is homologically
trivial on every smooth fibre of $T$. It defines a holomorphic
normal function $\nu$ in the family of intermediate Jacobians.
Its holomorphicity follows from the same locally defined chain
integrals: antiholomorphic derivatives vanish by the types of the
forms and of the algebraic relative cycle, and changing the chain
adds a period. The nonempty Zariski open $T$ in affine space is
analytically connected. For each $m>0$, the smooth fibre just
constructed has non-torsion value, so $m\nu$ cannot vanish identically
on any nonempty analytic open subset of $T$, even if that special
fibre lies outside $T^\circ$. Given any nonempty Zariski open subset
of $T$, choose one coordinate ball in its intersection with
$T^\circ$. On this ball the countably many period kernels and the
zero loci of $m\nu$, $m>0$, are all proper closed analytic subsets.
Baire's theorem shows that their joint complement is nonempty.
At each point in that complement, $m(L_0-L_1)$ is not algebraically
trivial for any $m>0$.

Suppose now that $m(L_0-L_1)$ were algebraically trivial on the
geometric generic fibre over $\overline K$. The assertion is
witnessed by finitely many curves, marked points, cycles and rational
equivalences. They descend to a finite extension of $K$. Normalize
the parameter space in this extension and shrink: the parameter
curves become smooth proper connected families with their sections,
and all finite cycle identities hold on every fibre. The image
contains a nonempty Zariski open subset of $T$. On its complex
points our cycle would be algebraically trivial, contradicting the
preceding paragraph. This proves \eqref{eq:specified-quintic-lines}.
Finally, algebraic equivalence does not acquire new relations under
an extension of algebraically closed fields: descend any proposed
finite relation to a finitely generated algebra over the smaller
field, spread its data over a nonempty open, and specialize at a
closed point, whose residue field is the smaller algebraically
closed field. This proves the assertion after every embedding into
$\mathbb C$.

In particular this proof supplies a nonzero rational Griffiths group
without importing Clemens's stronger infinite-dimensional theorem.
It also makes the counterexamples of C.1f completely specified:
start its induction with this $V$. The source has $G^i=0$, whereas
the target contains the displayed nonzero Griffiths class. Hard
Lefschetz implies that the image of the source meets the target
Griffiths group only in zero, so surjectivity fails for every ample
class. Infinite dimensionality is unnecessary for that conclusion.
The exact monomial count, both normal-derivative matrices, the
interpolation identities and the change-of-coordinates determinant
are recorded in \texttt{check\_quintic\_lines.py} and its output.
